\section{Model}
\label{sec:model}

This section presents the theoretical framework of the Agent Monte Carlo model. We establish the learning dynamics, prove equilibrium existence and stability, and derive comparative statics results.

\subsection{Environment}

Time is discrete, $t = 0, 1, \ldots, T$. There is a single risky asset with fixed supply normalized to 1. The fundamental value $V_t$ follows a geometric Brownian motion:
\begin{equation}
V_t = V_{t-1} \cdot \exp(\mu + \sigma_V \cdot \varepsilon_t), \quad \varepsilon_t \sim \mathcal{N}(0,1)
\end{equation}

\subsection{Agents}

There are $N$ agents in the economy. Each agent chooses from $K=5$ strategies:
\begin{itemize}
    \item \textbf{Noise Trader}: Random trading, provides liquidity
    \item \textbf{Momentum Trader}: Trends chasing based on past returns
    \item \textbf{Value Trader}: Mean-reversion based on fundamental value
    \item \textbf{Herd Trader}: Imitates other agents' behavior
    \item \textbf{Arbitrageur}: Exploits mispricing subject to constraints
\end{itemize}

Let $w_t \in \Delta(S)$ denote the strategy distribution at time $t$, where $\Delta(S) = \{w \in \mathbb{R}^K : w_k \geq 0, \sum_k w_k = 1\}$ is the simplex.

\subsection{Experience-Weighted Attraction Learning}

Agents learn using the Experience-Weighted Attraction (EWA) model of \citet{camerer1999experience}. The attraction $A_{k,t}$ of strategy $k$ evolves according to:
\begin{equation}
A_{k,t+1} = \frac{\phi \cdot N_t \cdot A_{k,t} + (\delta + (1-\delta) \cdot \mathbb{I}(s_k, s_t^*)) \cdot \pi_{k,t}}{N_{t+1}}
\end{equation}
where $N_{t+1} = \rho \cdot N_t + 1$ is the experience weight, $\phi \in [0,1]$ is the attraction decay rate, $\delta \in [0,1]$ is the imagination weight, and $\pi_{k,t}$ is the payoff of strategy $k$.

Strategy selection follows a logit rule:
\begin{equation}
P(s_{t} = s_k) = \frac{\exp(\lambda \cdot A_{k,t})}{\sum_{j=1}^K \exp(\lambda \cdot A_{j,t})}
\end{equation}
where $\lambda > 0$ is the selection sensitivity.

\subsection{Equilibrium Analysis}

\begin{definition}[EWA Equilibrium]
A strategy distribution $w^* \in \Delta(S)$ is an EWA equilibrium if:
\begin{equation}
w^*_k = \frac{\exp(\lambda A^*_k)}{\sum_{j=1}^K \exp(\lambda A^*_j)} \quad \forall k
\end{equation}
where $A^*$ satisfies the steady-state equation:
\begin{equation}
A^*_k = \frac{(\delta + (1-\delta) w^*_k) \cdot \pi_k(w^*)}{1 - \phi + (1-\rho)/N^*}
\end{equation}
with $N^* = 1/(1-\rho)$.
\end{definition}

\begin{assumption}
\label{ass:existence}
\begin{enumerate}
    \item[A1.] The strategy space $S$ is finite ($K < \infty$)
    \item[A2.] The payoff function $\pi_k(w)$ is continuous in $w$
    \item[A3.] Parameters satisfy: $\phi \in [0,1)$, $\delta \in [0,1]$, $\rho \in [0,1)$, $\lambda \in (0,\infty)$
\end{enumerate}
\end{assumption}

\begin{proposition}[Equilibrium Existence]
\label{prop:existence}
Under Assumption \ref{ass:existence}, there exists at least one EWA equilibrium $w^* \in \Delta(S)$.
\end{proposition}

\begin{proof}
Define the mapping $T: \Delta(S) \to \Delta(S)$ by:
\begin{equation}
T(w)_k = \frac{\exp(\lambda A_k(w))}{\sum_{j=1}^K \exp(\lambda A_j(w))}
\end{equation}
where $A_k(w) = \frac{(\delta + (1-\delta) w_k) \cdot \pi_k(w)}{C}$ with $C = 1 - \phi + (1-\rho)/N^* > 0$.

The simplex $\Delta(S)$ is compact and convex. The mapping $T$ is continuous because:
\begin{itemize}
    \item $\pi_k(w)$ is continuous by A2
    \item $A_k(w)$ is continuous (rational function with positive denominator)
    \item $\exp(\cdot)$ and the logit normalization are continuous
\end{itemize}

By Brouwer's fixed-point theorem, there exists $w^* \in \Delta(S)$ such that $T(w^*) = w^*$. This $w^*$ is an EWA equilibrium. $\square$
\end{proof}

\begin{assumption}
\label{ass:stability}
The Jacobian matrix $J = \frac{\partial T}{\partial w}\big|_{w^*}$ has spectral radius $\rho(J) < 1$.
\end{assumption}

\begin{proposition}[Local Stability]
\label{prop:stability}
Under Assumptions \ref{ass:existence} and \ref{ass:stability}, the equilibrium $w^*$ is locally asymptotically stable.
\end{proposition}

\begin{proof}
By the Hartman-Grobman theorem, the nonlinear system $w_{t+1} = T(w_t)$ can be linearized around $w^*$ as:
\begin{equation}
w_{t+1} - w^* \approx J \cdot (w_t - w^*)
\end{equation}
If $\rho(J) < 1$, the linear system is asymptotically stable, implying local asymptotic stability of the nonlinear system. $\square$
\end{proof}

\textbf{Remark on Global Stability:} Proposition 2 establishes local stability only. Global stability would require additional assumptions (e.g., contraction mapping). Numerically, we observe convergence from 20 random initial points to the same equilibrium, suggesting global stability in practice, though we leave formal proof for future research.

\subsection{Comparative Statics}

\begin{proposition}[Comparative Statics Derivative]
\label{prop:comparative}
The equilibrium strategy distribution $w^*(\theta)$ is differentiable with respect to parameter $\theta$, and:
\begin{equation}
\frac{\partial w^*}{\partial \theta} = \left[I - \frac{\partial T}{\partial w}\bigg|_{w^*}\right]^{-1} \cdot \frac{\partial T}{\partial \theta}
\end{equation}
\end{proposition}

\begin{proof}
Differentiating $w^* = T(w^*, \theta)$ with respect to $\theta$:
\begin{equation}
\frac{\partial w^*}{\partial \theta} = \frac{\partial T}{\partial w} \cdot \frac{\partial w^*}{\partial \theta} + \frac{\partial T}{\partial \theta}
\end{equation}
Rearranging and using the invertibility of $[I - \frac{\partial T}{\partial w}]$ (guaranteed by $\rho(J) < 1$):
\begin{equation}
\frac{\partial w^*}{\partial \theta} = \left[I - \frac{\partial T}{\partial w}\bigg|_{w^*}\right]^{-1} \cdot \frac{\partial T}{\partial \theta} \quad \square
\end{equation}
\end{proof}

\subsection{Welfare Analysis}

Individual utility for agent type $k$ is given by CARA utility with risk penalty:
\begin{equation}
U_k = -\exp(-\gamma \cdot W_k) - \frac{1}{2} \cdot \text{Var}(W_k)
\end{equation}
where $W_k$ is terminal wealth and $\gamma > 0$ is the risk aversion coefficient.

Social welfare follows a weighted utilitarian formulation:
\begin{equation}
\mathcal{W} = \sum_{k=1}^K \omega_k \cdot U_k
\end{equation}
where $\omega_k \geq 0$ are welfare weights summing to 1.

Inequality is measured by the Gini coefficient:
\begin{equation}
G = \frac{\sum_{i=1}^n \sum_{j=1}^n |W_i - W_j|}{2 n^2 \cdot \bar{W}}
\end{equation}
and the 90/10 wealth ratio $R_{90/10} = P_{90}(W) / P_{10}(W)$.

\subsection{Parameter Calibration}

Baseline parameters are calibrated from experimental literature:
\begin{itemize}
    \item $\phi = 0.88$, $\delta = 0.58$, $\rho = 0.85$, $\lambda = 1.5$ from \citet{camerer1999experience} Table III
    \item Payoff parameters from \citet{hvidkjaer2006trade} for momentum traders
    \item Herding parameters from \citet{shiller1995conversation}
\end{itemize}

\subsection{Numerical Verification}

We verify the theoretical results numerically:
\begin{itemize}
    \item Equilibrium existence: 100\% convergence rate (20/20 initial points)
    \item Residual: $3.2 \times 10^{-9}$ (numerical precision achieved)
    \item Spectral radius: $\rho(J) = 0.7342 < 1$ (locally stable)
    \item Parameter robustness: All 5 parameter sets tested are stable
\end{itemize}

\textbf{Equilibrium Uniqueness:} While we do not prove uniqueness analytically, our numerical experiments provide strong evidence. All 20 random initial points converge to the same equilibrium (pairwise correlation > 0.999). This suggests a unique equilibrium under baseline parameters, though multiple equilibria cannot be ruled out theoretically.

\begin{figure}[ht]
\centering
\includegraphics[width=\textwidth]{figures/equilibrium_convergence.pdf}
\caption{Equilibrium convergence and stability verification. Left: Convergence from 20 random initial points. Right: Jacobian eigenvalues (all inside unit circle).}
\label{fig:equilibrium}
\end{figure}
